% 第8章 功率放大电路
\chapter{功率放大电路}

\section{功率放大电路概述}
\subsection{功率放大电路的特点}
以输出功率为核心，要求：输出功率大、效率高、非线性失真小、晶体管安全工作。

\subsection{主要技术指标}
\begin{itemize}
    \item 输出功率 $P_O$：负载上获得的最大不失真功率；
    \item 效率 $\eta$：$\eta = \frac{P_O}{P_V} \times 100\%$（$P_V$ 为直流电源提供的功率）；
    \item 管耗 $P_T$：晶体管消耗的功率。
\end{itemize}

\subsection{功率放大电路的组成}
采用互补对称晶体管（NPN+PNP），工作在甲乙类/乙类状态，消除交越失真。

\section{互补功率放大电路}
\subsection{OCL电路的组成及工作原理}
OCL（无输出电容）电路采用双电源供电，互补管交替导通，负载获得完整正弦波，甲乙类工作方式消除交越失真。

\subsection{OCL电路的输出功率及效率}
最大输出功率：
\[ P_{OM} = \frac{V_{CC}^2}{2R_L} \]
效率（乙类）：
\[ \eta_{max} = \frac{\pi}{4} \approx 78.5\% \]

\subsection{OCL电路中晶体管的选择}
\subsubsection{最大管压降}
\[ |U_{CE_{max}}|=2V_{CC} \]

\subsubsection{集电极最大电流}
\[ I_{C_{max}}=\frac{V_{CC}}{R_L} \]

\subsubsection{集电极最大功耗}
\[ P_{T_{max}}=\frac{V^2_{CC}}{\pi^2R_L} \]
当 $U_{CES}=0$ 时：
\[ P_{T_{max}}=\frac{2}{\pi^2}P_{OM}\approx0.2P_{OM} \]

\section{功率放大电路的安全运行}
晶体管需满足：$U_{(BR)CEO} > 2V_{CC}$、$I_{CM} > \frac{V_{CC}}{R_L}$、$P_{CM} > 0.2P_{OM}$，必要时加散热片。

\section{集成功率放大电路}
集成运放型功率放大电路（如LM386），具有体积小、调试简单、性能稳定的特点，适用于小功率音频放大。

\section{常见功率放大电路}
包括OTL（有输出电容）、BTL（桥式推挽）等，OTL采用单电源供电，需输出电容隔直。

\section{例题参考}
4、9、11